\section{Problem Statement}

Regulatory knowledge is inherently interconnected, with legal documents frequently referencing other regulations, directives, standards, and technical specifications. These references form complex networks of dependencies in which understanding a specific legal requirement often requires traversing relationships between multiple documents. Previous work has highlighted that legal knowledge can be represented as interconnected networks or graphs, where documents and their references form nodes and edges that capture regulatory relationships.

A major challenge in establishing aforementioned graphs, is that the same regulation may be referenced using different surface forms throughout different legal texts. References may use the official title, legal identifier, common abbreviations, or descriptive names. For example, the regulation formally known as \textit{Regulation (EU) 2016/679} may be referenced as \textit{Regulation (EU) 2016/679}, \textit{General Data Protection Regulation}, or \textit{GDPR} depending on the context. Similarly, references to amending acts introduce additional ambiguity, as legal documents may refer to both the original regulation and subsequent amendments. For example, a legal text may reference \textit{Regulation (EU) 2016/679 (GDPR)} while another refers to a later regulation that modifies specific provisions of the GDPR. These cases create difficulties for automated retrieval systems, which must distinguish between the original legal entity, amending legislation, consolidated versions, and related documents.

Traditional keyword-based retrieval systems struggle with these challenges because they primarily rely on lexical overlap and often cannot distinguish between closely related legal entities. A search for a regulation identifier or title may retrieve the original regulation, an amending regulation, a consolidated version, or related documents without understanding the relationship between them. As a result, retrieving relevant regulatory knowledge requires methods capable of recognizing semantic equivalence between references and linking them to the correct underlying entities.

Therefore, a key prerequisite for enabling effective traversal and retrieval within regulatory knowledge graphs is reliable entity linking of legal references to specific regulatory entities. This thesis addresses this challenge by investigating AI-based approaches for regulatory entity linking and retrieval, enabling more accurate representation and exploration of interconnected regulatory knowledge.
